\section{Perspectiva del sistema: la división del trabajo entre Manager, Worker y Loop}
Entender el Agentic Loop de veRL no consiste solo en memorizar unos cuantos nombres de clases, sino en ver con claridad cómo separa el razonamiento de múltiples turnos, la invocación de herramientas y la gestión de servicios. El Manager se encarga de la orchestration, el Worker se encarga de la ejecución y LoopBase se encarga de definir la interfaz abstracta; esta división en capas permite que el sistema reutilice la infraestructura entre distintas tareas.

\begin{knowledgebox}{Por qué esta arquitectura conviene al Agentic RL}
\begin{itemize}
  \item Un bucle de múltiples turnos requiere de manera natural una gestión de estado unificada
  \item La invocación de herramientas y el razonamiento del modelo por lo general deben reutilizarse en módulos separados
  \item Durante el entrenamiento hay que atender al mismo tiempo el rollout, el server y la lógica de la tarea
\end{itemize}
\end{knowledgebox}

\subsection{Resumen de la sección}
Lo importante del Agentic Loop no es qué tan compleja sea alguna clase en particular, sino que abstrae la lógica de interacción think-act-observe en un esqueleto de sistema reutilizable.

\newpage
